% Generated by build/sync_qft_soln.py from committed sources only.
% Source repository: https://github.com/Luca-Yucheng-Jin/QFT-soln
% Source commit: 07d7fe95fd5cf5b26d38a16ee535a04a925277b2
\section{Exact Results (PSI QFT II, Tutorial 4)}

\noindent\textit{Question prompts paraphrased from the official
\href{https://psi-online.perimeterinstitute.ca/courses/take/quantum-field-theory-ii-student/pdfs/20776379-module-4-tutorial-1}{Perimeter PSI Online sheet};
the equations and notation follow the source.}

\subsection{The K\"all\'en--Lehmann Representation}

For a free scalar field, recall the Feynman propagator
\begin{equation}
    D_0^F(x-y)
    \equiv \langle 0|T[\phi(x)\phi(y)]|0\rangle
    =-i\int\frac{d^4k}{(2\pi)^4}\,
    \frac{e^{ik\cdot(x-y)}}{k^2+m_0^2}.
\end{equation}
It describes propagation of a particle of mass $m_0$. Its pole at
$k^2=-m_0^2$ is understood with the usual $-i\epsilon$ prescription. In an
interacting theory, replace the free vacuum by the interacting vacuum and define
\begin{equation}
    D^F(x-y)=\langle\Omega|T[\phi(x)\phi(y)]|\Omega\rangle.
\end{equation}
The lowest pole of this propagator determines the physical mass. To obtain its
nonperturbative spectral representation, insert a complete set of interacting
states between the two fields.

\begin{enumerate}[label=\alph*)]
    \item Assume Poincar\'e invariance, with four-momentum operator
    $P^\mu=(H,\mathbf P)$ satisfying
    \begin{equation}
        [P^\mu,P^\nu]=0.
    \end{equation}
    Determine the action of translations on the field and the interacting vacuum:
    \begin{equation}
        e^{-iP\cdot x}\phi(0)e^{iP\cdot x}=?
        \qquad
        e^{-iP\cdot x}|\Omega\rangle=?
    \end{equation}
    \begin{solution}
    The vacuum should be invariant under Poincaré transformation, namely $e^{-iP\cdot x}\ket{\Omega} = \ket{\Omega}$. It can also be shown that
    \begin{equation}
        e^{-iP\cdot x}\phi(0)e^{iP\cdot x}= \phi(x).
    \end{equation}
    \end{solution}
    Take the states to be simultaneous energy--momentum eigenstates. Denote a
    one-particle state of physical mass $m$ by $|\mathbf k\rangle$, with
    $E_{\mathbf k}=\sqrt{\mathbf k^2+m^2}$, and use
    $|\mathbf k,n\rangle$ for states containing two or more particles, where
    $n$ collects the remaining quantum numbers.

    \item Write the spectral decomposition of the identity, including the vacuum,
    one-particle states, and multiparticle states.
    \begin{solution}
        \begin{equation}
            1 = \ket{\Omega}\bra{\Omega} + \int\frac{d^3k}{(2\pi)^32E_\mathbf{k}}\ket{\mathbf{k}}\bra{\mathbf{k}} + \sum_n\int\frac{d^3k}{(2\pi)^32E_\mathbf{k}}\ket{\mathbf{k},n}\bra{\mathbf{k},n}.
        \end{equation}
    \end{solution}

    \item For $x^0>y^0$, insert that identity on the right-hand side of the
    interacting propagator and show that
    \begin{align}
        \langle\Omega|\phi(x)\phi(y)|\Omega\rangle
        &={}
        \int\frac{d^3\mathbf k}{(2\pi)^3}\,
        \frac{e^{ik\cdot(x-y)}}{2E_{\mathbf k}}
        \left|\langle\mathbf k|\phi(0)|\Omega\rangle\right|^2
        \notag\\
        &\quad+
        \sum_n\int\frac{d^3\mathbf k}{(2\pi)^3}\,
        \frac{e^{ik\cdot(x-y)}}{2E_{\mathbf k}}
        \left|\langle\mathbf k,n|\phi(0)|\Omega\rangle\right|^2.
    \end{align}
    Assume that $\phi(x)|\Omega\rangle$ has a one-particle contribution.
    \begin{solution}
        \begin{equation}
            \begin{aligned}
                 \langle\Omega|\phi(x)\phi(y)|\Omega\rangle &=  \langle\Omega|\phi(0)e^{iP\cdot (x-y)}\phi(0)|\Omega\rangle\\
                 &=\langle\Omega|\phi(0)e^{iP\cdot (x-y)}\bigg(\ket{\Omega}\bra{\Omega} + \int\frac{d^3k}{(2\pi)^32E_\mathbf{k}}\ket{\mathbf{k}}\bra{\mathbf{k}} \\
                 &\qquad+ \sum_n\int\frac{d^3k}{(2\pi)^32E_\mathbf{k}}\ket{\mathbf{k},n}\bra{\mathbf{k},n}\bigg)\phi(0)|\Omega\rangle\\
                 &=\int\frac{d^3\mathbf k}{(2\pi)^3}\,
        \frac{e^{ik\cdot(x-y)}}{2E_{\mathbf k}}
        \left|\langle\mathbf k|\phi(0)|\Omega\rangle\right|^2
        +
        \sum_n\int\frac{d^3\mathbf k}{(2\pi)^3}\,
        \frac{e^{ik\cdot(x-y)}}{2E_{\mathbf k}}
        \left|\langle\mathbf k,n|\phi(0)|\Omega\rangle\right|^2.
            \end{aligned}
        \end{equation}
    \end{solution}

    \item Introduce the spectral function
    \begin{equation}
        \rho(s)
        \equiv
        \delta(s-m^2)
        \left|\langle\mathbf k|\phi(0)|\Omega\rangle\right|^2
        +\sum_n
        \left|\langle\mathbf k,n|\phi(0)|\Omega\rangle\right|^2
        \delta(s-M_n^2),
        \qquad
        M_n^2=E_{\mathbf k}^2-\mathbf k^2.
    \end{equation}
    Assuming there are no bound states, show that $\rho(s)=0$ for $s<4m^2$
    except at $s=m^2$, whereas $\rho(s)$ is nonzero for $s\geq4m^2$. Use
    this result to establish, for $x^0>y^0$,
    \begin{align}
        \langle\Omega|\phi(x)\phi(y)|\Omega\rangle
        &={}
        \int\frac{d^3\mathbf k}{(2\pi)^3}\,
        \frac{e^{ik\cdot(x-y)}}{2E_{\mathbf k}}
        \left|\langle\mathbf k|\phi(0)|\Omega\rangle\right|^2
        \notag\\
        &\quad+
        \int_{4m^2}^{\infty}ds\,\rho(s)
        \int\frac{d^3\mathbf k}{(2\pi)^3}\,
        \frac{e^{ik\cdot(x-y)}}{2E_{\mathbf k}},
    \end{align}
    and, for $x^0<y^0$,
    \begin{align}
        \langle\Omega|\phi(y)\phi(x)|\Omega\rangle
        &={}
        \int\frac{d^3\mathbf k}{(2\pi)^3}\,
        \frac{e^{-ik\cdot(x-y)}}{2E_{\mathbf k}}
        \left|\langle\mathbf k|\phi(0)|\Omega\rangle\right|^2
        \notag\\
        &\quad+
        \int_{4m^2}^{\infty}ds\,\rho(s)
        \int\frac{d^3\mathbf k}{(2\pi)^3}\,
        \frac{e^{-ik\cdot(x-y)}}{2E_{\mathbf k}}.
    \end{align}
    \begin{solution}
        If there are no bound states, then it is obvious that there are no additional poles for $s<4m^2$. The result follows trivially.
    \end{solution}

    \item Use the QFT I identity
    \begin{align}
        \int\frac{d^4k}{(2\pi)^4}\,
        \frac{e^{ik\cdot(x-y)}}{k^2+m^2-i\epsilon}
        &={}
        i\Theta(x^0-y^0)
        \int\frac{d^3\mathbf k}{(2\pi)^3}\,
        \frac{e^{ik\cdot(x-y)}}{2E_{\mathbf k}}
        \notag\\
        &\quad+
        i\Theta(y^0-x^0)
        \int\frac{d^3\mathbf k}{(2\pi)^3}\,
        \frac{e^{-ik\cdot(x-y)}}{2E_{\mathbf k}}
        \label{eq:contourint}
    \end{align}
    to derive the K\"all\'en--Lehmann representation
    \begin{equation}
        \widetilde D_F(k^2)
        =\frac{-i\left|\langle\mathbf k|\phi(0)|\Omega\rangle\right|^2}
               {k^2+m^2-i\epsilon}
        +\int_{4m^2}^{\infty}ds\,\rho(s)
        \frac{-i}{k^2+s-i\epsilon},
    \end{equation}
    where $\widetilde D_F$ is the Fourier transform of $D_F(x)$.
    \begin{solution}
        This follows trivially from repeating the same contour integral one used to obtain \eqref{eq:contourint}.
    \end{solution}

    \item Regard $\widetilde D_F(k^2)$ as an analytic function of $k^2$. Locate
    its lowest pole and explain the physical meaning of both the pole and its residue.
    \begin{solution}
        The lowest pole is at $k^2 = -m^2$, which is the physical mass of the one-particle state, and the residue is the field renormalisation factor $Z$.
    \end{solution}

    \item The continuum begins at $s=4m^2$ and produces a branch cut. If the
    theory also contains bound states, which determine where their contributions would
    appear in $\rho(s)$.
    \begin{solution}
        The bound states would still be single poles,m they would appear as $\delta$-functions in the spectral function.
    \end{solution}
\end{enumerate}

\subsection{Schwinger--Dyson Equations}

The Schwinger--Dyson equations state that correlation functions obey the
classical equations of motion up to local contact terms.

\begin{enumerate}[label=\alph*)]
    \item In the free Euclidean Klein--Gordon theory, use invariance of the path
    integral under $\phi(x)\mapsto\phi(x)+\epsilon(x)$. Starting from
    \begin{equation}
        \langle\phi(x)\rangle
        =\left.
        \frac{1}{\mathcal Z[0]}
        \frac{\delta\mathcal Z[J]}{\delta J(x)}
        \right|_{J=0},
    \end{equation}
    show that
    \begin{equation}
        (-\partial_z^2+m^2)
        \langle\phi(z)\phi(x)\rangle
        =\hbar\delta(z-x).
    \end{equation}
    \begin{solution}
        Consider the integral under an infinitesimal field shifting $\phi \to \phi +\epsilon$
        \begin{equation}
            \int\mathcal{D}\phi F[\phi]e^{-S[\phi]/\hbar} = \int \mathcal{D}\phi F[\phi+\epsilon]e^{-S[\phi+\epsilon]/\hbar}.
        \end{equation}
        Then, we have
        \begin{equation}
            \begin{aligned}
                0 &=\int \mathcal{D}\phi\left[\left(F[\phi] + \frac{\delta F}{\delta\phi}\epsilon\right)\left(1- \frac{1}{\hbar}\frac{\delta S}{\delta\phi}\epsilon\right) - F[\phi]\right]e^{-S[\phi]/\hbar}\\
                &=\left\langle\frac{\delta F[\phi]}{\delta\phi }\right\rangle\epsilon  -\frac{1}{\hbar} \left\langle F[\phi]\frac{\delta S[\phi]}{\delta \phi}\right\rangle\epsilon
            \end{aligned}
        \end{equation}
        for all $\epsilon$. Therefore, we find
        \begin{equation}
            \hbar \left\langle\frac{\delta F[\phi]}{\delta\phi }\right\rangle = \left\langle F[\phi]\frac{\delta S[\phi]}{\delta \phi}\right\rangle.
        \end{equation}
        For our case, we take
        \begin{equation}
            S_0[\phi] = \int d^dz \frac{1}{2}\phi(z)(-\partial_z^2 +m^2)\phi(z), \quad F[\phi] = \phi(x),
        \end{equation}
        which gives
        \begin{equation}
            (-\partial_z^2+m^2)
        \langle\phi(z)\phi(x)\rangle
        =\hbar\delta(z-x).
        \end{equation}
    \end{solution}

    \item Add the interaction
    \begin{equation}
        S[\phi]
        =S_0[\phi]+\frac{\lambda}{n!}\int d^4y\,\phi^n(y),
    \end{equation}
    and show that
    \begin{equation}
        (-\partial_z^2+m^2)
        \langle\phi(z)\phi(x)\rangle
        =-
        \left\langle
            \frac{\lambda}{(n-1)!}\phi^{n-1}(z)\phi(x)
        \right\rangle
        +\hbar\delta(z-x).
    \end{equation}
    \begin{solution}
        Now, we have
        \begin{equation}
            \frac{\delta S[\phi]}{\delta\phi} = (-\partial_z^2+m^2)\phi(z) + \frac{\lambda}{(n-1)!}\phi^{n-1}(z),
        \end{equation}
        and hence, using the master equation, we find
        \begin{equation}
                    (-\partial_z^2+m^2)
        \langle\phi(z)\phi(x)\rangle
        =-
        \left\langle
            \frac{\lambda}{(n-1)!}\phi^{n-1}(z)\phi(x)
        \right\rangle
        +\hbar\delta(z-x).
        \end{equation}
    \end{solution}
    \item Use invariance of the generating functional under
    $\phi\mapsto\phi+\delta\phi$ to prove the general identity
    \begin{equation}
        \left\langle
            \frac{\delta S}{\delta\phi(x)}
            \prod_{r=1}^{n}\phi(x_r)
        \right\rangle
        =\hbar\sum_{j=1}^{n}
        \delta^4(x-x_j)
        \left\langle
            \prod_{\substack{r=1\\r\neq j}}^{n}\phi(x_r)
        \right\rangle.
    \end{equation}
    What changes in real time?
    \begin{solution}
        This follows directly from substituting
        \begin{equation}
            F[\phi] = \prod_{r=1}^n\phi(x_r).
        \end{equation}
        In real time, the expression becomes
        \begin{equation}
             \left\langle
            \frac{\delta S}{\delta\phi(x)}
            \prod_{r=1}^{n}\phi(x_r)
        \right\rangle
        =i\hbar\sum_{j=1}^{n}
        \delta^4(x-x_j)
        \left\langle
            \prod_{\substack{r=1\\r\neq j}}^{n}\phi(x_r)
        \right\rangle.
        \end{equation}
    \end{solution}

    \item Show that
    \begin{equation}
        \langle\phi_1\cdots\phi_n\rangle
        =\int d^4x\,
        \Delta_{x1}(-\partial_x^2+m^2)
        \langle\phi_x\phi_2\cdots\phi_n\rangle,
    \end{equation}
    where $\phi_i=\phi(x_i)$ and
    $\Delta_{ij}=\Delta_F(x_i-x_j)$.
    \begin{solution}
        Using $(-\partial_x^2+m^2)\Delta_{1x} = \delta_{1x}$, we obtain the above equation.
    \end{solution}

    \item In the free Euclidean Klein--Gordon theory, combine the preceding result
    with the Schwinger--Dyson equations to obtain
    \begin{equation}
        \langle\phi_1\cdots\phi_4\rangle_0
        =\Delta_{12}\Delta_{34}
        +\Delta_{13}\Delta_{24}
        +\Delta_{14}\Delta_{23}.
    \end{equation}
    \begin{solution}
        \begin{equation}
            \begin{aligned}
                \langle\phi_1\phi_2\phi_3\phi_4\rangle_0 &= \int d^dx \Delta_{x1}(-\partial_x^2 +m^2)\langle\phi_x\phi_2\phi_3\phi_4\rangle_0\\
                &=\int d^dx \Delta_{x1}(\delta_{x2}\Delta_{34} + \delta_{x3}\Delta_{24} + \delta_{x4}\Delta_{23})\\
                &=\Delta_{12}\Delta_{34}
        +\Delta_{13}\Delta_{24}
        +\Delta_{14}\Delta_{23},
            \end{aligned}
        \end{equation}
        where we will take $\hbar =1$ from now on.
    \end{solution}

    \item In Euclidean $\phi^3$ theory, apply the same identities recursively and
    show that
    \begin{align}
        \langle\phi_1\phi_2\rangle
        &={}
        \Delta_{12}
        +(-g)^2\int d^4x\,d^4y\,
        \left(
            \frac{1}{2}\Delta_{x1}\Delta_{xy}^2\Delta_{y2}
            +\frac{1}{4}\Delta_{x1}\Delta_{xx}\Delta_{yy}\Delta_{y2}
        \right)
        +O(g^4).
    \end{align}
    Does this agree with the result obtained in Tutorial 2?
    \begin{solution}
        \begin{equation}
            \begin{aligned}
                \langle\phi_1\phi_2\rangle &= \int d^dx \Delta_{1x}(-\partial_x^2 + m^2)\langle\phi_x\phi_2\rangle\\
                &=  \int d^dx \Delta_{1x}\left(-\frac{g}{2}\langle\phi_x^2\phi_2\rangle + \delta   _{x2}\right)\\
                &= \Delta_{12} - \frac{g}{2}\int d^dx d^dy \Delta_{1x}\Delta_{2y}(-\partial_y^2 + m^2)\langle\phi_x^2\phi_y\rangle\\
                &= \Delta_{12} - \frac{g}{2}\int d^dx d^dy \Delta_{1x}\Delta_{2y}\left(-\frac{g}{2}\langle\phi^2_y\phi_x^2\rangle + 2\delta_{xy}\langle\phi_x\rangle\right).
            \end{aligned}
        \end{equation}
        Since we are only interested to order $g^2$, we have
        \begin{equation}
            \langle\phi^2_y\phi_x^2\rangle  = \langle\phi^2_y\phi_x^2\rangle_0 + \mathcal{O}(g^2) = 2\Delta_{xy}^2 + \Delta_{xx}\Delta_{yy}+ \mathcal{O}(g^2).
        \end{equation}
        Meanwhile, the tadpole diagram is given by
        \begin{equation}
        \begin{aligned}
            \langle\phi_x\rangle &= \int d^dz \Delta_{xz}(-\partial_z^2+m^2)\langle\phi_z\rangle\\
            &=-\frac{g}{2}\int d^dz \Delta_{xz}\langle\phi_z^2\rangle\\
            &=-\frac{g}{2}\int d^dz \Delta_{xz}\Delta_{zz} + \mathcal{O}(g^3).
        \end{aligned}
        \end{equation}
        Then, we have
        \begin{equation}
            \begin{aligned}
                \langle\phi_1\phi_2\rangle &= \Delta_{12} +(-g)^2\int d^dx d^dy \Delta_{1x}\Delta_{2y}\left(\frac{1}{2}\Delta_{xy}^2 + \frac14\Delta_{xx}\Delta_{yy}\right)\\
                &\qquad+ (-g)^2\int d^dx d^dy\,\frac{1}{2}\Delta_{1x}\Delta_{2x}\Delta_{xy}\Delta_{yy} + \mathcal{O}(g^4).
            \end{aligned}
        \end{equation}
    \end{solution}

    \item \textbf{Optional.} For a continuous symmetry of the Lagrangian, prove
    the Ward--Takahashi identity
    \begin{align}
        \partial_\mu
        \langle j^\mu(x)\phi(x_1)\cdots\phi(x_n)\rangle
        ={}&
        \hbar\sum_{j=1}^{n}
        \bigl\langle
            \phi(x_1)\cdots\phi(x_{j-1})
            \notag\\[-0.4em]
            &\qquad\times
            \delta\phi(x)\delta^4(x-x_j)
            \phi(x_{j+1})\cdots\phi(x_n)
        \bigr\rangle,
    \end{align}
    where, in the notation of the sheet,
    \begin{equation}
        j^\mu(x)
        =\frac{\delta S}{\delta\phi_{,\mu}(x)}\delta\phi(x)
        -\delta\mathcal L(x),
        \qquad
        \phi_{,\mu}\equiv\partial_\mu\phi.
    \end{equation}
    What changes in real time?
    \begin{solution}
        The classical action has some global symmetry, $\phi(x) \to \phi(x) + \epsilon(x) \delta\phi(x)$, which corresponds to a Noether current, $j_\mu(x)$.
        Now, we promote the symmetry to be a local symmetry $\phi(x) \to \phi(x) + \epsilon(x)\delta\phi(x)$. Again, we consider the path integral
        \begin{equation}
            \int \mathcal{D}\phi F[\phi] e^{-S[\phi]/\hbar} = \int  \mathcal{D}\phi (F[\phi]+\delta F[\phi])\left(1 -  \frac{1}{\hbar}\delta S[\phi]\right)e^{-S[\phi]/\hbar},
        \end{equation}
        where $F[\phi] \to F[\phi] + \delta F[\phi]$ under the symmetry transformation. Note that $\delta S, \delta F$ are of order $\epsilon$, and for an infinitesimal transformation, we can ignore terms of order $\epsilon^2$. Then, we have
        \begin{equation}
            \hbar \langle\delta F[\phi]\rangle = \langle\delta S[\phi]F[\phi]\rangle.
        \end{equation}
        Recall
\begin{equation}
    \begin{aligned}
        \delta S
        &=
        \int d^d x\,
        \left[
            \frac{\partial\mathcal L}{\partial\phi(x)}
            \epsilon(x)\delta\phi(x)
            +
            \frac{\partial\mathcal L}
            {\partial(\partial_\mu\phi(x))}
            \partial_\mu\!\left(
                \epsilon(x)\delta\phi(x)
            \right)
        \right]\\
        &=
        \int d^d x\,\epsilon(x)
        \left[
            \frac{\partial\mathcal L}{\partial\phi(x)}
            \delta\phi(x)
            +
            \frac{\partial\mathcal L}
            {\partial(\partial_\mu\phi(x))}
            \partial_\mu\delta\phi(x)
        \right] +
        \int d^d x\,
        \frac{\partial\mathcal L}
        {\partial(\partial_\mu\phi(x))}
        \delta\phi(x)\partial_\mu\epsilon(x)\\
        &=
        \int d^d x\,
        \left[
            \epsilon(x)\partial_\mu
            \delta\mathcal L^\mu(x)
            +
            \frac{\partial\mathcal L}
            {\partial(\partial_\mu\phi(x))}
            \delta\phi(x)\partial_\mu\epsilon(x)
        \right]\\
        &=
        \int d^d x\,
        \partial_\mu\!\left(
            \epsilon(x)\delta\mathcal L^\mu(x)
        \right)
        +
        \int d^d x\,
        \left[
            \frac{\partial\mathcal L}
            {\partial(\partial_\mu\phi(x))}
            \delta\phi(x)
            -
            \delta\mathcal L^\mu(x)
        \right]
        \partial_\mu\epsilon(x)\\
        &=
        \int d^d x\,
        j^\mu(x)\partial_\mu\epsilon(x)\\
        &=
        -\int d^d x\,
        \epsilon(x)\partial_\mu j^\mu(x).
    \end{aligned}
\end{equation}
        and take
        \begin{equation}
        \begin{aligned}
                        F[\phi] = \phi(x_1) \cdots \phi(x_n) \implies \delta F[\phi] &=\sum_{j=1}^{n}
            \phi(x_1)\cdots\phi(x_{j-1})
             \epsilon(x_j)\delta\phi(x_j)
            \phi(x_{j+1})\cdots\phi(x_n)\\
            &=\sum_{j=1}^{n}\int d^dx \delta^d(x-x_j)
            \phi(x_1)\cdots\phi(x_{j-1})
             \epsilon(x)\delta\phi(x)
            \phi(x_{j+1})\cdots\phi(x_n).
        \end{aligned}
        \end{equation}

        Hence, we find the Ward-Takahashi identity
        \begin{align}
        \partial_\mu
        \langle j^\mu(x)\phi(x_1)\cdots\phi(x_n)\rangle
        ={}&
        -\hbar\sum_{j=1}^{n}
        \bigl\langle
            \phi(x_1)\cdots\phi(x_{j-1})
            \delta\phi(x)\delta^4(x-x_j)
            \phi(x_{j+1})\cdots\phi(x_n)
        \bigr\rangle.
        \end{align}
        In real time, this becomes
                \begin{align}
        \partial_\mu
        \langle j^\mu(x)\phi(x_1)\cdots\phi(x_n)\rangle
        ={}&
        -i\hbar\sum_{j=1}^{n}
        \bigl\langle
            \phi(x_1)\cdots\phi(x_{j-1})
            \delta\phi(x)\delta^4(x-x_j)
            \phi(x_{j+1})\cdots\phi(x_n)
        \bigr\rangle.
        \end{align}
    \end{solution}

    \item \textbf{More optional.} Show that contact terms on the right-hand side
    of the Schwinger--Dyson and Ward--Takahashi identities do not contribute to
    $i\mathcal M$. In Euclidean momentum conventions, use
    \begin{equation}
        \langle f|S|i\rangle
        =\lim_{k_i^2\to-m^2}
        (k_i^2+m^2)\cdots
        \langle\Omega|T\widetilde\phi(k_1)\cdots|\Omega\rangle,
    \end{equation}
    with
    \begin{equation}
        \widetilde\phi(k)
        =i\int d^4x\,e^{-ik\cdot x}\phi(x).
    \end{equation}
    The scattering amplitude and momentum-space correlator contain the same
    momentum-conserving delta function:
    \begin{equation}
        \langle f|S|i\rangle
        =(2\pi)^4\delta^4\!\left(\sum k\right)i\mathcal M,
    \end{equation}
    \begin{equation}
        \langle\Omega|T\widetilde\phi(k_1)\cdots|\Omega\rangle
        =(2\pi)^4\delta^4\!\left(\sum k\right)\mathcal F(k_i).
    \end{equation}
    Near the simultaneous poles,
    \begin{equation}
        \mathcal F(k_i)
        =\frac{i\mathcal M}
        {(k_1^2+m^2)\cdots(k_n^2+m^2)}
        +\text{nonsingular terms}.
    \end{equation}
    Argue that the Fourier transform of a contact term lacks the required
    multivariable pole and hence cannot contribute to $i\mathcal M$.
    \begin{solution}
        Consider a contact term
        \begin{equation}
            \delta^d(x_1-x_2)\langle\phi(x_3) \cdots \phi(x_n)\rangle,
        \end{equation}
        which is Fourier transformed to
        \begin{equation}
            (2\pi)^d\delta^d(k_1 +k_2)\langle\tilde \phi(k_3) \cdots \tilde \phi(k_n)\rangle.
        \end{equation}
        Its contribution to the LSZ is
        \begin{equation}
            \lim_{k_i^2\to-m^2}(k_1^2+m^2)\cdots(k_n^2+m^2)(2\pi)^d\delta^d(k_1 +k_2)\langle\tilde \phi(k_3) \cdots \tilde \phi(k_n)\rangle.
        \end{equation}
        Since it contains no poles in $k_1,k_2$, it has no contribution to the LSZ reduction formula.
    \end{solution}

    \item \textbf{More optional.} In Lorenz gauge, use the photon LSZ formula
    \begin{equation}
        \langle f|S|i\rangle
        =\epsilon^\mu\int d^4x\,e^{-ik\cdot x}
        (-\partial^2)\cdots
        \langle\Omega|T A_\mu(x)\cdots|\Omega\rangle
    \end{equation}
    together with the Ward--Takahashi identity to derive
    \begin{equation}
        k_\mu\mathcal M^\mu=0,
        \qquad
        \mathcal M=\epsilon_\mu\mathcal M^\mu.
    \end{equation}
    \begin{solution}
        In QED, we have the equation of motion for the gauge fields in the Lorenz gauge
        \begin{equation}
            -\partial^2A_\mu = j_\mu = e\bar \psi \gamma_\mu \psi.
        \end{equation}
        The current corresponds to the global $U(1)$ transformation. The LSZ becomes
        \begin{equation}
            \langle f|S|i\rangle=\epsilon^\mu\int d^dx\,e^{-ik\cdot x}
        \langle\Omega|T j_\mu(x)\cdots|\Omega\rangle = (2\pi)^d\delta^d\left(\sum k\right)\epsilon_\mu \mathcal{M}_\mu.
        \end{equation}
        Now replacing $\epsilon^\mu \to k^\mu$, we have
        \begin{equation}
        \begin{aligned}
            (2\pi)^d\delta^d\left(\sum k\right)k_\mu \mathcal{M}_\mu &= i\int d^dx  (\partial^\mu e^{-ik\cdot x}) \langle\Omega|T j_\mu(x)\cdots|\Omega\rangle\\
            &= - i\int d^dx   e^{-ik\cdot x}\partial^\mu\langle\Omega|T j_\mu(x)\cdots|\Omega\rangle.
        \end{aligned}
        \end{equation}
        Now, since we have shown that the contact terms do not contribute to $S$-matrix, and the LHS of the Ward-Takahashi identity is purely contact terms. Therefore, we can conclude that $k_\mu \mathcal{M}^\mu = 0$.
    \end{solution}
\end{enumerate}
